\section{Derivations}
\label{app:derivations}

% Mirrors Wu & Tu Appendix A, which carries the inference details the main text compresses
% (their A.5 is where the message weights \lambda_Z, \lambda_H are justified).

\subsection{Full mean-field updates}
\label{app:full-updates}

Fix a slot $t$. The prefix is frozen as $\qbar{1},\dots,\qbar{t-1}$ and contracted into the
arc--label scores $\Bfac{c}_{j,a}$ of \Cref{eq:contract}, with
$\Bfac{c}_{\textsc{root},a} = \rfac{c}_a$; the live variables are $\QW{t}$ on $\Vocab$,
$\QZ{t}$ on the $\nlab$ labels, and one $\QH{t}{c}$ on $\Dep{t}$ per channel. Taking the
negative log of the step conditional \Cref{eq:unclamp} in expectation under the product
$\QW{t}\,\QZ{t}\prod_c \QH{t}{c}$ gives the per-slot energy
%
\begin{equation}
\label{eq:slot-energy}
\begin{aligned}
E_t &= -\textstyle\sum_w \QW{t}(w)\,\bfac_w\\
    &\quad - \textstyle\sum_{w,a} \QW{t}(w)\,\QZ{t}(a)\,\Sfac_{w,a}\\
    &\quad - \textstyle\sum_c \sum_{j\in\Dep{t}} \sum_a
             \QH{t}{c}(j)\,\QZ{t}(a)\,\Bfac{c}_{j,a},
\end{aligned}
\end{equation}
%
the local normaliser having contributed a constant that no derivative sees. Two features of
\Cref{eq:slot-energy} do the work below. It is multilinear, so the partial derivative with
respect to any one variable does not contain that variable and the coordinate step is available
in closed form. And it is one energy per slot rather than one energy for the sentence, which is
how \Cref{thm:trilemma} is respected rather than contradicted: the prefix enters $E_t$ only as
the constants $\Bfac{c}_{j,a}$, while the factors in which slot $t$ serves a later slot $u>t$
belong to $E_u$, where $\qbar{t}$ is in turn a constant and not a variable of the problem. No
term is deleted from $E_t$ by hand; the anti-causal one is never formed, which is the
distinction \Cref{rmk:never-formed} draws.

The three gradients are, with nothing removed,
%
\begin{align}
\label{eq:grad-W}
  \frac{\partial E_t}{\partial \QW{t}(w)}
    &= -\bfac_w - \textstyle\sum_a \QZ{t}(a)\,\Sfac_{w,a},\\
\label{eq:grad-H}
  \frac{\partial E_t}{\partial \QH{t}{c}(j)}
    &= -\textstyle\sum_a \QZ{t}(a)\,\Bfac{c}_{j,a},\\
\label{eq:grad-Z}
  \frac{\partial E_t}{\partial \QZ{t}(a)}
    &= -\textstyle\sum_w \QW{t}(w)\,\Sfac_{w,a}\\
    &\quad - \textstyle\sum_c \sum_{j\in\Dep{t}} \QH{t}{c}(j)\,\Bfac{c}_{j,a}.\nonumber
\end{align}
%
What inference minimises is not $E_t$ itself but $E_t$ regularised by the entropies of its three
variables, each of them carrying its own weight,
%
\begin{equation}
\label{eq:slot-free-energy}
\begin{aligned}
  \Free_t = E_t &- \lamW \Ent(\QW{t}) - \lamZ \Ent(\QZ{t})\\
                &- \lamH \textstyle\sum_c \Ent\big(\QH{t}{c}\big),
\end{aligned}
\end{equation}
%
which is the extended entropic Frank--Wolfe objective of \citet{wu2023probabilistic}, their
Appendix~A.2, and is the ordinary mean-field free energy when all three weights equal one. Their
Appendix~A.3 solves the resulting subproblem in closed form:
$\min_{z\in\Delta}\langle c,z\rangle + \lambda\sum_s z_s\log z_s$ is attained at
$z = \softmax(-c/\lambda)$. Applying it coordinate by coordinate to
\Cref{eq:grad-W,eq:grad-H,eq:grad-Z} gives the three updates
%
\begin{align}
\label{eq:full-update-H}
  \QH{t}{c}(j) &\propto \exp\Big(\tfrac{1}{\lamH}
      \textstyle\sum_a \QZ{t}(a)\,\Bfac{c}_{j,a}\Big),\\
\label{eq:full-update-Z}
  \QZ{t}(a) &\propto \exp\Big(\tfrac{1}{\lamZ}\Big[
      \textstyle\sum_w \QW{t}(w)\,\Sfac_{w,a}\\
    &\qquad\quad + \textstyle\sum_c\sum_{j\in\Dep{t}}
      \QH{t}{c}(j)\,\Bfac{c}_{j,a}\Big]\Big),\nonumber\\
\label{eq:full-update-W}
  \QW{t}(w) &\propto \exp\Big(\tfrac{1}{\lamW}\Big[
      \bfac_w + \textstyle\sum_a \QZ{t}(a)\,\Sfac_{w,a}\Big]\Big).
\end{align}
%
Because $\softmax$ is unchanged by adding a constant to every coordinate, only the centred
gradient --- its projection onto the tangent space of the simplex --- is ever read, which is why
the constant discarded from \Cref{eq:slot-energy} is immaterial.
\Cref{eq:full-update-H,eq:full-update-Z} are \Cref{eq:update-H,eq:update-Z} of the main text with
the word message $\sum_w \QW{t}(w)\Sfac_{w,a}$ in place of the observed unary;
\Cref{eq:full-update-W} is the word update, and its softmax over $\Vocab$ is the language-model
output layer.

\paragraph{Schedule.}
Each variable is initialised from its own unary-type factors, the rule of
\citet{wu2023probabilistic}'s Equation~7: $\QW{t}^{(0)} \propto \exp(\bfac)$ and
$\QZ{t}^{(0)} \propto \exp(\tilde{s}/\lamZ)$ with
$\tilde{s}_a = \sum_w \QW{t}^{(0)}(w)\Sfac_{w,a}$. The asynchronous order of their \S2.3.2 then
applies verbatim --- update $\{\QH{t}{c}\}$, then $\QZ{t}$ --- for $\tau$ rounds, finishing with
one application of \Cref{eq:full-update-W} as the readout. Taking $\tau \ge 2$ is what makes the
attention query context-dependent; at $\tau=1$ it is the fixed learned probe
$\softmax(\tilde{s}/\lamZ)$, which is legal but strictly weaker. Iterating the
$W\!\leftrightarrow\!Z$ loop beyond the readout is also legitimate mean field, but it feeds the
model's own word guess back into the label belief, so we take the readout before feedback as the
default and treat extra word rounds as an ablation.

\paragraph{The observed mode is the same energy.}
Clamping $\QW{t}$ to the point mass $\delta_{w_t}$ turns the word message
$\sum_w \QW{t}(w)\Sfac_{w,a}$ into $\Sfac_{w_t,a}$, makes \Cref{eq:full-update-W} idle, and
reduces \Cref{eq:slot-energy} to
%
\begin{equation}
\label{eq:obs-energy}
\begin{aligned}
  E_t^{\text{obs}} &= -\textstyle\sum_a \QZ{t}(a)\,\Sfac_{w_t,a}\\
    &\quad -\textstyle\sum_c \sum_{j\in\Dep{t}} \sum_a
      \QH{t}{c}(j)\,\QZ{t}(a)\,\Bfac{c}_{j,a},
\end{aligned}
\end{equation}
%
up to the additive constant $-\bfac_{w_t}$, which no gradient sees. \Cref{eq:obs-energy} is
exactly the per-position energy of the conditional chain of \Cref{sec:chain}, and its two
gradients are exactly \Cref{eq:update-H,eq:update-Z}. Encoding and decoding are therefore one
energy family under two conditioning patterns, which is the formal content of
\Cref{fig:factor-graph} and the reason $\Sfac$ is a single factor trained in both of its roles.
A second consequence is visible in \Cref{eq:grad-H}: the head gradient contains no word message
at all, so the attention update is the same functional in both modes and the two runs differ
only through the belief at which it is evaluated. Sharing one attention pattern between the
predictive and the observed pass at a slot is therefore not an approximation imposed from
outside but a choice of initialisation, which mean field leaves free.

\subsection{Derivation of the exact readout}
\label{app:exact-derivation}

With the prefix frozen, the factors of \Cref{eq:unclamp} that remain live at slot $t$ are the
word unary $\bfac$, the word--label factor $\Sfac$, and one contracted arc factor per channel.
Their graph is a star centred on $\Zv{t}$, with leaves $\Wv{t}$ and
$\Hv{t}{1},\dots,\Hv{t}{\nchan}$; a star is a tree, so sum--product terminates after one sweep
and returns the exact marginal. We eliminate the leaves in turn.

Each head variable appears in exactly one term of the exponent and carries no unary of its own,
so its elimination is a plain sum over $\Dep{t}$ under counting measure, and the sums over
distinct channels separate:
%
\begin{equation}
\label{eq:elim-H}
\begin{aligned}
  \sum_{j_1,\dots,j_\nchan \in \Dep{t}}
    \exp\Big(\textstyle\sum_c \Bfac{c}_{j_c,a}\Big)
  &= \prod_c \sum_{j\in\Dep{t}} e^{\Bfac{c}_{j,a}}\\
  &=: \readout{t}(a).
\end{aligned}
\end{equation}
%
The message from channel $c$ into $\Zv{t}$ is thus $\sum_{j\in\Dep{t}} e^{\Bfac{c}_{j,a}}$, in
logarithms $\LSE_{j\in\Dep{t}} \Bfac{c}_{j,a}$, seeded at $j = \textsc{root}$ by $\rfac{c}_a$,
and $\log\readout{t}(a)$ is the sum of these $\nchan$ messages. Eliminating $\Zv{t}$ against the
word--label factor then leaves
%
\begin{equation}
\label{eq:elim-Z}
\begin{aligned}
  \hat{p}\big(\Wv{t}{=}w \mid w_{\prefix{t}}\big)
    &\;\propto\; \textstyle\sum_a e^{\bfac_w + \Sfac_{w,a}}\,\readout{t}(a)\\
    &\;=\; e^{\bfac_w} \textstyle\sum_a e^{\Sfac_{w,a}}\,\readout{t}(a),
\end{aligned}
\end{equation}
%
which is \Cref{eq:exact-readout}. The slot normaliser cancels in the softmax over $\Vocab$, so
the implementation works in logits, $\bfac_w + \LSE_a\,(\Sfac_{w,a} + \log\readout{t}(a))$, and
never forms $\readout{t}$ itself.

Writing $\Znorm_a = \sum_{w} e^{\bfac_{w} + \Sfac_{w,a}}$ and
$p(w \mid a) = e^{\bfac_w + \Sfac_{w,a}}/\Znorm_a$ puts \Cref{eq:elim-Z} in the form
$\hat p(\cdot) = \sum_a \pi_t(a)\, p(\cdot \mid a)$ with mixture weights
$\pi_t(a) \propto \readout{t}(a)\,\Znorm_a$. The $\nlab$ components are fixed by the parameters
and only the weights depend on the context, which is the precise sense in which
\Cref{prop:readout} calls this a mixture of softmaxes, and the precise sense in which it is a
weaker family than one whose components are themselves context-dependent.

\paragraph{The sequence form.}
Along the sequence the head domain grows by exactly one position per step,
$\Dep{t+1} = \Dep{t}\cup\{t\}$, so the channel message of \Cref{eq:elim-H} is a causal prefix
log-sum-exp in $t$. Evaluated slot by slot it would cost $O(\seqlen^2\nlab\nchan)$; evaluated as
a prefix reduction it is one \texttt{logcumsumexp} scan along the position axis,
$O(\seqlen\nlab\nchan)$, for all $t$ at once.

The clipped distance table refines this without changing it. The arc score of an arc spanning
distance $t-j$ is drawn from bucket $f(t-j)$ of \citet{wu2023probabilistic}'s clip function, of
which the causal model keeps only the $t-j>0$ half: distances $1,\dots,\clip$ each own a bucket
and every distance beyond $\clip$ shares the last, giving $\clip+1$ buckets. Writing
$\Bfac{c,k}_{j,a} = \sum_b \qbar{j}(b)\,\Tfac{c,k}_{a,b}$ for the contraction against bucket
$k$'s own arc matrix,
%
\begin{equation}
\label{eq:mu-scan}
\begin{aligned}
  \log\readout{t}^{(c)}(a) = \LSE\Big(\;&\rfac{c}_a,\;\;
      \LSE_{1\le j\le t-\clip-1} \Bfac{c,\clip+1}_{j,a},\\[-2pt]
    &\Bfac{c,1}_{t-1,a},\;\dots,\;\Bfac{c,\clip}_{t-\clip,a}\Big),
\end{aligned}
\end{equation}
%
where $\log\readout{t}(a) = \sum_c \log\readout{t}^{(c)}(a)$ and any term whose position index
falls below $1$ is $-\infty$. Only the far bucket accumulates: its positions form a contiguous
prefix, and it is the one \texttt{logcumsumexp}, offset by $\clip+1$. Each of the $\clip$ near
buckets contains exactly one position at every slot, so it is a plain shift of that bucket's
table by its own distance, and the readout is one log-sum-exp over $\clip+2$ aligned tensors.
The global variable of \Cref{app:globals}, being one more leaf off $\Zv{t}$, contributes one
further term $\LSE_m \Bglob_{m,a}$ to $\log \readout{t}(a)$ and leaves this argument untouched.

\paragraph{The layer-parallel schedule.}
Three things have to hold for it to survive, and all three do. The readout is closed-form on top of the final
prefix beliefs, so it does not enter the iteration and cannot serialise it. The scan reads only
$j<t$ and is a prefix reduction, hence associative and computable for all slots in one pass
rather than by a loop over $t$. And the query stream disappears entirely: with
\Cref{eq:exact-readout} as the readout there is no predictive inner loop and no second set of
attention problems, so a layer remains one triangular-masked pass exactly as in
\Cref{sec:variants}. At generation time the same structure gives the readout a cache of its own
beside the content stream's, namely a running $\nlab \times \nchan$ table of channel messages,
extended by a single \texttt{logaddexp} for each token that is generated.

\subsection{Message weights}
\label{app:message-weights}

The weights $\lamZ,\lamH,\lamW,\lamG$ are the per-variable regularisation weights of
\Cref{eq:slot-free-energy}. They are not free rescalings of a fixed objective: each is the
entropy weight of its own variable, so changing one changes which quantity is being minimised,
and a sweep over $\lamH$ is a sweep over a family of objectives rather than over a temperature
within one.

\paragraph{$\lamZ$ and $\lamH$.}
\citet{wu2023probabilistic} fix $\lamZ = 1$ and $\lamH = 1/\nlab$, and their Appendix~A.5
derives the second from a variance argument in the manner of \citet{vaswani2017attention}.
Assume the arc scores are independent with mean zero and variance $\sigma^2$. The head logit of
\Cref{eq:full-update-H} is the double contraction
$\sum_{a,b} \QZ{t}(a)\,\qbar{j}(b)\,\Tfac{c}_{a,b}$; at uniform beliefs each of its $\nlab^2$
terms carries weight $\nlab^{-2}$, so the logit has variance $\sigma^2/\nlab^2$, and dividing by
$\lamH = 1/\nlab$ restores $\sigma^2$. Nothing in that computation is specific to the encoder:
the causal head logit has the same shape, with one of the two beliefs frozen rather than live,
so $\lamH = 1/\nlab$ transfers unchanged. Their argument for $\lamZ$ is negative --- the
corresponding variance for the label message depends on the sentence length, which varies, so no
fixed $\lamZ$ recovers $\sigma^2$, and they set it to one for simplicity. Two things change in
the causal model and neither rescues it. The label message here sums over $\Dep{t}$, whose size
is $t$, so the quantity that obstructed a fixed $\lamZ$ now varies from position to position
within a sentence and not only between sentences; and the factor of two carried by the encoder's
label message, which comes from symmetrising over head--child order, is absent, since the
directed model traverses each arc once. We make the same choice and set $\lamZ = 1$.
\NOTE{the position-dependence of $|\Dep{t}|$ is our transposition of their variance argument to
the causal head domain; it is not something \citet{wu2023probabilistic} state.}

\paragraph{$\lamW$.}
The word weight has no counterpart in \citet{wu2023probabilistic}, whose encoder has no output
layer, and it has none in the mainline readout either: \Cref{eq:elim-Z} is untempered
sum--product, in which $\Wv{t}$ is eliminated exactly rather than variationally, so no variable
is left to carry a weight. The question arises only in the mean-field ablation, and there two
arguments pull against each other; both are worth stating. The first fixes $\lamW = 1$ so that
the readout is an untempered conditional, on the ground that with $\bfac$ and $\Sfac$ learned,
any other value is absorbed into them. The second reopens the choice, and is the one we find
decisive: absorption is available only if $\Sfac$ serves one role, whereas weight tying makes
the scale of $\Sfac$ the output logit range \emph{and} the strength of the unary at once, so
rescaling $\Sfac$ to absorb $\lamW$ rescales the unary with it. A weight $\lamW<1$ is then the
one in-model lever that decouples output sharpness from unary scale, and it is a legitimate
per-variable message weight in exactly the sense $\lamH$ already is. The trade is worth stating
rather than hiding: at $\lamW \ne 1$ the readout is a tempered conditional, and the claim that
the model is calibrated by construction weakens to calibrated up to one declared temperature. We
report $\lamW = 1$ throughout unless stated otherwise.
\NOTE{ours, and drawn nowhere else --- a consequence for \Cref{sec:exp3}: the evaluation-time swap
between the two readouts is only meaningful at $\lamW = 1$, since at any other value the
mean-field readout is tempered and the exact one is not, so the swap would measure the
temperature along with the approximation.}

\paragraph{$\lamG$.}
Nothing fixes the weight of the global variable. \citet{wu2023probabilistic} introduce it in
their Appendix~B.3 without a message weight and never run it, and the construction that carries
it into the causal model (\Cref{app:globals}) gives its update but not its temperature. We set
$\lamG = \lamZ = 1$.
This is a declared choice and not a derivation: $\Gv{t}$ is a leaf off $\Zv{t}$ whose update is
a softmax over the $\nglob$ global components, no variance argument of the kind that fixes
$\lamH$ is available for it, and we had no reason to prefer another value. It is not swept in
any run reported here. \NOTE{ours, and untuned --- worth a line in \Cref{sec:exp-ablations} if
the global-variable arm turns out to be sensitive to it.}

\paragraph{Calibration and the uniform-belief limit.}
All of the above calibrates $\lamH$ at one point of the model's trajectory. The variance
argument evaluates the head logit at uniform beliefs, which is where the model starts and not
where it stays; \Cref{sec:temperature} works out what that costs, why the error grows with the
label width, and what the belief-normalised variants are. We do not repeat that analysis here
and only record the connection: $\lamH = 1/\nlab$ is the right constant for the uniform limit of
the derivation above, and every variant of \Cref{sec:temperature} is built to agree with it
exactly at that limit.

\subsection{Complexity}
\label{app:complexity}

\begin{table}[t]
\centering
\small
\begin{tabular}{@{}lcc@{}}
\toprule
Layer & Width term & Length term \\
\midrule
Transformer layer          & $12\dmodel^2$ & $2\seqlen\dmodel$ \\
PT content layer           & $3\nlab^2$    & $2\seqlen\nlab$   \\
PT, mean-field two-stream  & $5\nlab^2$    & $4\seqlen\nlab$   \\
PT, exact-readout mainline & $3\nlab^2$    & $2\seqlen\nlab$   \\
\bottomrule
\end{tabular}
\caption{Approximate FLOPs per token and per layer, under the calibration
$\krank\nchan \approx \nlab$ for the arc-score decomposition and at the matched width
$\nlab = \dmodel$ that \Cref{sec:cmp-readout} argues is the right comparison unit. The
transformer's $12\dmodel^2$ is $4\dmodel^2$ of query, key, value and output projections plus
$8\dmodel^2$ of feed-forward block; the causal PT's $3\nlab^2$ is $3\nlab\krank\nchan$ --- the
prefix keys formed once, since keys are values, plus the query contraction and the output
contraction --- and it has no feed-forward term at all. The two-stream row shares its keys
between the streams. The exact-readout mainline runs the content layer alone and adds, per token
rather than per layer, $O(\nlab\nchan)$ for the $\readout{t}$ scan and $O(|\Vocab|\nlab)$ for
the vocabulary readout.}
\label{tab:complexity}
\end{table}

\Cref{tab:complexity} accounts for one layer at one position. The headline is not the one the
construction invites: the causal PT is not twice a transformer layer but roughly half of one,
because the feed-forward block that dominates a transformer layer is absent while the attention
term is the same in both. Even the mean-field two-stream variant, which does pay for a second
set of inference problems, sits at roughly $0.4$--$0.6\times$ a transformer layer, and the
exact-readout mainline removes the second stream altogether. Adding the global variables of
\Cref{app:globals} costs a further $\approx 2\nglob\nlab$ per token, which narrows the gap
without closing it.

Two costs sit outside the table because they are per token rather than per layer. Generation
extends the readout's $\nlab\times\nchan$ cache by one \texttt{logaddexp}, $O(\nlab\nchan)$,
which is cheaper than any attention pass and is the whole of what the readout adds to the
incremental step. The vocabulary readout itself is $O(|\Vocab|\nlab)$, the cost every language
model pays for its output layer, but it is a log-sum-exp rather than a matrix product: the same
arithmetic with worse hardware constants, since the exponentials are element-wise and the
reduction cannot be handed to a matrix unit. The materialised $|\Vocab|\times\nlab$ intermediate
is the real hazard, and the implementation chunks the reduction over the vocabulary in the
manner of a fused cross-entropy, accumulating logits blockwise so that the intermediate never
exists in full.

Two qualifications keep the table honest. The first is that a per-layer FLOP count is half a
comparison when one of the models shares its parameters across iterations: at a matched
parameter budget the causal PT can afford more iterations per parameter than a transformer can
afford layers, so depth and parameters cannot both be matched and which one is matched has to be
declared rather than inferred. \Cref{sec:settings} matches the parameter budget across the three
models and sets the loop depth of the Looped baseline to the iteration count $\niter$, and
reports a sweep over $\niter$ so that the comparison does not rest on one point of that
trade.\NOTE{keep this in step with \Cref{sec:settings}, where the exact parameter-matching rule
--- total or non-embedding --- is still to be fixed.} The second is that none of this predicts
wall-clock. \citet{wu2023probabilistic} measure their encoder at about three times slower than a
transformer despite having far fewer parameters, and attribute it in their Appendix~C to the way
relative positions are encoded: the arc matrices are gathered from the clipped distance table
per position pair, which is parameter-inefficient and, more to the point here, memory-bound.
That axis is orthogonal to how many streams a model runs and to everything in
\Cref{tab:complexity}, and it is why \Cref{sec:settings} reports wall-clock beside the parameter
count instead of deriving one from the other.
